\scanpage{533}%
\mspage{444}%
% archive sheet number: 532
ment, supposons que $\underline{x}$ soit ess. constant, \deleted{\unsure{$\exists$|$\forall$}} \deleted{\unsure{$A$}} et \unsure{prouvons|prenons} \inserted{\unsure{que|pour}} \deleted{\unsure{le cri|la cat.}} \unsure{l'on|l'on y} peut réindexer comme indiqué$^{*}$.

% margin note written sideways
\margin{Le cas où $I$ admet un objet final est trivial, \unsure{on prendra|prendre} $J=e$, $p=\mathrm{id}_{x_i}$ Supposons donc \unsure{que $I$ n'ait pas d'objet final|qu'on n'ait pas d'objet final}.}

Comme $I$ grd devant $\pi_0$ et $\pi_0\geq\operatorname{card}\operatorname{Ob} C$, on voit qu'il existe un \inserted{$x_0\in\operatorname{Ob} C$ tel que la} sous-catégorie pleine $I'\inserted{=I_{x_0}}$ de $I$, \deleted{cofinale} formée des $i$ tels que $x_i=x_0$, soit cofinale ds $I$. \deleted{Donc} Comme $(x_i)_{i\in I}$ est ess. constant, il en est de $=$ de $(x_{i'},i'\in I')$, et on peut lui appliquer le critère 4.15.7.1, avec $J\to I'$ cofinal. La démonstration montre de plus que, si $I'$ est grd devant $\pi_0$ (\inserted{cardinal infini} sans autres conditions) OPS $J$ grde devant $\pi_0$. \inserted{De plus, il est clair qu'on peut supposer $J$ ordonné.}

% slanted
\margin{\unsure{en tenir compte dans la dém. de|en tenir compte ds la dém. de} 4.15.7.1.}

Comme $\pi_0\geq\operatorname{card}\operatorname{Hom}(x_0,x_0)$, on voit aussi qu'il existe \deleted{\unsure{un $p\in$|un $p$}} \inserted{un projecteur} $p$ ds $x_0$ tel que la partie \deleted{cofinale} $\inserted{J_p=}J'$ de $J$, \deleted{telle que} formée des $j$ tels que $p_j=p$, soit cofinale. Étant cofinale, elle est égalemt grde devant $\pi_0$. Donc OPS que tous les $p_j$ sont égaux à un même $p$. Les conditions $1^\circ$--$4^\circ$ deviennent

\[1^\circ)\quad p\,\tau_{j'j}=\tau_{j'j}\,p\quad\text{pour }j\leq j'.\]

\[2^\circ)\quad\deleted{p\tau_{j'j}=\tau_{j'j}}\quad\forall j,\ \text{existe }j'\geq j\ \text{avec}\quad p\tau_{j'j}=\tau_{j'j}p=\tau_{j'j}.\]

$3^\circ)$ \deleted{\unsure{$\forall$}} L'image \inserted{$x_0^\natural$} du projecteur $p$ dans $x_0$ existe.

$4^\circ)$ L'endom. $\tau_{j'j}^\natural$ de $x_0^\natural$ est iso, pour $j\leq j'$.
